\subsection{A quantitative extremal estimate and all large-boundary limits}

We retain the notation $s_i=(L_i/2\pi)^2$, $F_i(s)=\log r_i(s)$,
and $K(p)$ of Section~\ref{sec:A2b-secondary}. All radii use fixed tangent
coordinates on the fixed marked compactification $Y$.

\begin{lemma}[Proved: a uniform weighted-radius deficit]\label{A2-strebel-deficit}
Put $c=2g-2+n>0$. For every $s\in(0,\infty)^I$ and every pairwise
disjoint system of pointed conformal discs $E_i$ avoiding the other
marked points,
\[
 \sum_{i\in I}s_i\bigl(\log r_i(E_i)-F_i(s)\bigr)\leq c.
\]
Consequently
\[
 0\leq K(s)-s\cdot F(s)\leq c.
\]
The constants do not depend on the boundary lengths. Fixed cusp ends
are allowed.
\end{lemma}
\begin{proof}
Let $g_s$ be the metric which is hyperbolic on the bordered core and
flat on each appended cylinder. Its core has area $2\pi c$ by
Gauss--Bonnet, including the finite-area cusp ends. Every closed curve
in the peripheral free homotopy class of $x_i$ has $g_s$-length at least
$L_i$. Indeed, the map collapsing every flat cylinder orthogonally onto
its boundary is length nonincreasing. The resulting curve in the
hyperbolic core has length at least its geodesic representative, of
length $L_i$. This argument also covers a curve originally contained
entirely in the cylinder.

Let $\zeta_i:E_i\to\D$ be its centred uniformization. On the annulus
$\epsilon<|\zeta_i|<R$, $R<1$, write the metric as
$\rho_i^2|d\zeta_i|^2$. Each circle is peripheral, so
\[
 \int_0^{2\pi}\rho_i(re^{i\theta})r\,d\theta\geq L_i.
\]
Cauchy--Schwarz, followed by integration in $dr/r$, gives
\[
 \operatorname{Area}_{g_s}\{\epsilon<|\zeta_i|<R\}
 \geq\frac{L_i^2}{2\pi}\log(R/\epsilon).
\]
The test annuli are disjoint. Near $x_i$, the canonical cap coordinate
$w_i$ and the test coordinate satisfy
$w_i=(r_i(E_i)/r_i(s))\zeta_i+O(\zeta_i^2)$, up to a rotation.
The flat metric is $s_i|dw_i/w_i|^2$. Thus, after deleting the small
discs $|\zeta_i|\leq\epsilon$ at all active ends, the area of the whole
surface is
\[
 2\pi c+2\pi\sum_i s_i
 \bigl(F_i(s)-\log r_i(E_i)-\log\epsilon\bigr)+o(1).
\]
The error tends to zero for fixed $s$: the analytic coordinate change
has a nonzero first derivative, and its higher terms contribute
$O(\epsilon)$ to the boundary height. Compare this area with the sum
of the lower bounds, first let $\epsilon\downarrow0$, then let
$R\uparrow1$, and cancel the divergent terms. The first assertion
follows. Taking the supremum over the test systems gives the upper
bound in the second assertion by the weighted reduced-modulus
characterization of $K$. The lower bound holds because the canonical
caps themselves are a competing disc system.
\end{proof}

\begin{theorem}[Proved with the stated fixed-compactification hypothesis:
the cap limit]\label{A2-cap-limit}
Fix $Y$ and its marked points as above. Let $s^{(m)}>0$ and $t_m\to\infty$
satisfy $s^{(m)}/t_m\to p\in[0,\infty)^I\setminus\{0\}$.
For every $i$ with $p_i>0$, the pointed caps $D_i(s^{(m)})$ converge in
the Carath\'eodory kernel sense to the Jenkins--Strebel disc $U_i(p)$
with circumference weights $2\pi\sqrt{p_i}$. Their normalized conformal
maps converge locally uniformly, and
\[
 r_i(s^{(m)})\longrightarrow R_i(p).
\]
For $p_i=0$, the radii tend to zero. In particular the limit depends
on the relative growth of the lengths and on $Y$; it is not universally
a Koebe slit domain. No assertion of a fixed limit is made when the
marked compactification itself is allowed to diverge.
\end{theorem}
\begin{proof}
Every competing pointed disc has a finite uniform radius bound $B_i$.
For a sphere this is Lemma~\ref{A1-cap-bound} after normalizing two
other marked points to $1,\infty$. For a torus, lift to its flat
universal covering: injectivity downstairs bounds the area of the
lifted image by the torus area, and the area formula bounds its first
Taylor coefficient. For genus at least two, lift to the uniformizing
disc and use Schwarz's lemma. These arguments also give compactness
of normalized univalent maps whenever their radii stay bounded away
from zero. In the spherical case use the usual distortion bound for
normalized univalent functions; the bounded first coefficient makes
the family locally bounded. In the torus case the same distortion
bound applies to the lifts. In higher genus use normality of maps to
the uniformizing disc. A nonconstant limit remains injective, and
Hurwitz's theorem preserves omission of the other marked points.

Let $J=\{i:p_i>0\}$. Choose any fixed competing system $E_i$ for all
active points. Divide Lemma~\ref{A2-strebel-deficit} by $t_m$. Terms
with $i\notin J$ on the left side $\sum s_i^{(m)}F_i(s^{(m)})/t_m$
are bounded above by $\sum_{i\notin J}(s_i^{(m)}/t_m)\log B_i=o(1)$.
It follows that
\[
 \liminf_m\sum_{i\in J}\frac{s_i^{(m)}}{t_m}F_i(s^{(m)})
 \geq\sum_{i\in J}p_i\log r_i(E_i).
\]
The same estimate and the individual upper bounds imply a uniform
lower bound for each $F_i(s^{(m)})$, $i\in J$. Hence we may replace
$s_i^{(m)}/t_m$ by $p_i$ in this inequality. The supremum of the right
side is $K(p)$: approximate the positive-weight Strebel discs by compact
concentric subdiscs, and put sufficiently small discs at the remaining
points. Conversely, feasibility gives
$\sum_{i\in J}p_iF_i(s^{(m)})\leq K(p)$. Therefore this weighted sum
converges to $K(p)$.

From any subsequence extract convergent normalized cap maps at every
positive-weight point. Their nonconstant limits have disjoint images:
an overlap would put a common small open set into the images for all
large $m$, by local uniform convergence and the inverse mapping theorem.
Their radii are the limits of the first derivatives. Thus their disc
system attains the weighted reduced-modulus maximum $K(p)$.
The uniqueness clause in the Jenkins--Strebel extremal theorem makes
it precisely $\{U_i(p):i\in J\}$. This identifies every convergent
subsequence and proves the asserted convergence for the full sequence.
The locally uniform convergence of normalized univalent maps is
equivalent to the stated pointed kernel convergence.

Finally these positive-weight Strebel discs fill $Y$ outside their
critical graph, which has empty interior. If a zero-weight cap had
radius bounded away from zero along a subsequence, compactness would
give an additional nonempty open disc disjoint from all the limiting
positive-weight discs. Such an open set cannot lie in the critical
graph. This contradiction proves its radius tends to zero.
\end{proof}

\begin{corollary}[Known form, with explicit derivation: the Koebe and
orange limits]\label{A2-explicit-limits}
Normalize a capped pant by placing its centres at $0,1,\infty$.
If $L_0\to\infty$ and $L_1/L_0,L_\infty/L_0\to0$, then the cap at
$0$ tends to $\C\setminus[1,\infty)$, its conformal radius tends to
$4$, and its normalized conformal map tends to
\[
 k(u)=\frac{4u}{(1+u)^2},\qquad |u|<1.
\]
If instead $L_0/L_1,L_0/L_\infty\to1$ and the lengths tend to infinity,
the cap at $0$ tends to
\[
 Q_0^\circ=\{|z|<1,\ \operatorname{Re}z<1/2\},
 \qquad r(Q_0,0)=\frac4{3\sqrt3}.
\]
The other two limits follow by the permutations of $0,1,\infty$.
This is the precise limiting sense in which the orange slices are valid.
\end{corollary}
\begin{proof}
Up to the common positive scale, the three-pole Strebel quadratic
differential with squared circumference parameters $p_0,p_1,p_\infty$
has the rational expression
\[
 q_p=-\frac{p_\infty z^2+(p_1-p_\infty-p_0)z+p_0}
              {z^2(z-1)^2}\,dz^2.
\]
There is no accessory parameter: the difference between two such
quadratic differentials would have at most three simple poles on the
sphere and must vanish. For $p=(1,0,0)$ this is
$-dz^2/(z^2(1-z))$. Direct substitution of $z=k(u)$ gives
$k^*q_p=-du^2/u^2$. The map $k$ is conformal onto the stated slit
domain, as follows either by its explicit inverse using
$\sqrt{1-z}$ with value $1$ at $0$, or by its boundary image and
nonvanishing derivative in $\D$. Its derivative at zero is $4$.

For $p=(1,1,1)$ the numerator is $z^2-z+1$, with simple zeros
$a=e^{i\pi/3}$ and $b=e^{-i\pi/3}$. On the long unit-circle arc
joining $a$ to $b$, and on the vertical segment joining them,
$q_p$ is positive on the tangent directions. Together with the long
arc of $|z-1|=1$ joining the same vertices, these three symmetric
arcs form the critical graph. The complementary disc about $0$ is
$Q_0^\circ$. To check its radius explicitly, put
$m(z)=(z-a)/(z-b)$; it maps this lune to a sector of angle $2\pi/3$.
The branch $v(z)=(e^{-i\pi/3}m(z))^{3/2}$ maps that sector to the
upper half-plane and takes $0$ to $i$. Therefore
$U(z)=(v(z)-i)/(v(z)+i)$ is a centred disc uniformization.
Since $m'(0)/m(0)=i\sqrt3$, we obtain
$|U'(0)|=3\sqrt3/4$, giving the displayed radius. This also verifies
the cylinder structure of these complementary domains directly.
Theorem~\ref{A2-cap-limit} now proves both limits.
\end{proof}

\begin{corollary}[Known inequality with derivation: the image is not a box]
\label{A2-nonbox}
For the pants normalization $0,1,\infty$, every genuine admissible cap
system satisfies $r_0r_1<1$. Nevertheless the individual supremum of
each of $r_0,r_1,r_\infty$ over such systems is $4$. The set of
admissible radii is therefore not a Cartesian product of intervals.
\end{corollary}
\begin{proof}
For $p=(1,1,0)$ the quadratic differential is
$-dz^2/[z^2(z-1)^2]$. Its two disc domains are the half-planes separated
by $\operatorname{Re}z=1/2$; each has conformal radius $1$ at its marked
centre. The weighted reduced-modulus inequality gives $r_0r_1\leq1$,
the two-domain bound stated in the opened original
\cite[Part I, \S2, Theorem3, pp.~169--173]{ALavrentiev1934}:
his maps go from the domains to the unit disc, so their derivative
product bound is the reciprocal conformal-radius bound. It is strict for genuine
canonical caps: their disjoint analytic boundaries admit simultaneous
small enlargements into the nonempty hyperbolic core, which strictly
increase their radii. The individual suprema are at most $4$ by Koebe
and are approached by the one-dominant-boundary limits just proved.
A Cartesian product with these individual suprema would contain two
radii both larger than $1$, contradicting the product inequality.
\end{proof}
